\begin{ex}
Gốc của tia số biểu diễn số tự nhiên nào?
\shortans{0}
\loigiai{
Gốc tia số bắt đầu từ số 0.
}
\end{ex}

\begin{ex}
Trên tia số, điểm đứng liền sau số 14 là số nào?
\shortans{15}
\loigiai{
Số liền sau số 14 là 15.
}
\end{ex}

\begin{ex}
Trên tia số, điểm đứng liền trước số 28 là số nào?
\shortans{27}
\loigiai{
Số liền trước số 28 là 27.
}
\end{ex}

\begin{ex}
Tìm số liền sau của số 69.
\shortans{70}
\loigiai{
Số liền sau số 69 là 70.
}
\end{ex}

\begin{ex}
Tìm số liền trước của số 90.
\shortans{89}
\loigiai{
Số liền trước số 90 là 89.
}
\end{ex}

\begin{ex}
Tìm số liền trước của số liền trước số 35.
\shortans{33}
\loigiai{
Số liền trước 35 là 34, liền trước 34 là 33.
}
\end{ex}

\begin{ex}
Tìm số liền sau của số liền sau số 48.
\shortans{50}
\loigiai{
Số liền sau 48 là 49, liền sau 49 là 50.
}
\end{ex}

\begin{ex}
Trên tia số, số nào nằm chính giữa số 30 và số 32?
\shortans{31}
\loigiai{
Số nằm chính giữa là 31.
}
\end{ex}

\begin{ex}
Trên tia số, khoảng cách từ vạch số 5 đến vạch số 12 là bao nhiêu đơn vị?
\shortans{7}
\loigiai{
Khoảng cách: $12 - 5 = 7$ đơn vị.
}
\end{ex}

\begin{ex}
Trên tia số, khoảng cách từ vạch số 20 đến vạch số 35 là bao nhiêu đơn vị?
\shortans{15}
\loigiai{
Khoảng cách: $35 - 20 = 15$ đơn vị.
}
\end{ex}

\begin{ex}
Điểm A trên tia số ở vạch 8, điểm B ở vạch 14. Hỏi từ điểm A phải bước thêm bao nhiêu bước đơn vị để đến điểm B?
\shortans{6}
\loigiai{
Số bước: $14 - 8 = 6$ bước.
}
\end{ex}

\begin{ex}
Một chú thỏ đang ở vạch số 15 trên tia số, thỏ nhảy tiến về phía trước 4 bước. Chú thỏ dừng lại ở vạch số mấy?
\shortans{19}
\loigiai{
Dừng lại ở: $15 + 4 = 19$.
}
\end{ex}

\begin{ex}
Một chú ếch đang ở vạch số 24 trên tia số, ếch nhảy lùi lại 6 bước. Chú ếch dừng lại ở vạch số mấy?
\shortans{18}
\loigiai{
Dừng lại ở: $24 - 6 = 18$.
}
\end{ex}

\begin{ex}
Tìm tổng của số liền trước số 20 và số liền sau số 20.
\shortans{40}
\loigiai{
Số liền trước là 19, liền sau là 21. Tổng: $19 + 21 = 40$.
}
\end{ex}

\begin{ex}
Tìm hiệu giữa số liền sau số 50 và số liền trước số 50.
\shortans{2}
\loigiai{
Số liền sau là 51, liền trước là 49. Hiệu: $51 - 49 = 2$.
}
\end{ex}

\begin{ex}
Số liền sau của số lớn nhất có hai chữ số là số nào?
\shortans{100}
\loigiai{
Số lớn nhất có hai chữ số là 99, số liền sau là 100.
}
\end{ex}

\begin{ex}
Số liền trước của số bé nhất có hai chữ số là số nào?
\shortans{9}
\loigiai{
Số bé nhất có hai chữ số là 10, số liền trước là 9.
}
\end{ex}

\begin{ex}
Trên tia số, giữa số 42 và số 47 có bao nhiêu số tự nhiên?
\shortans{4}
\loigiai{
Các số ở giữa: 43, 44, 45, 46. Có 4 số.
}
\end{ex}

\begin{ex}
Trên tia số từ vạch 0 đến vạch 10 có bao nhiêu đoạn đơn vị?
\shortans{10}
\loigiai{
Có đúng 10 đoạn đơn vị.
}
\end{ex}

\begin{ex}
Điểm M biểu diễn số 50. Điểm N nằm bên phải điểm M và cách M 8 đơn vị. Điểm N biểu diễn số nào?
\shortans{58}
\loigiai{
Điểm N biểu diễn: $50 + 8 = 58$.
}
\end{ex}

\begin{ex}
Điểm P biểu diễn số 75. Điểm Q nằm bên trái điểm P và cách P 15 đơn vị. Điểm Q biểu diễn số nào?
\shortans{60}
\loigiai{
Điểm Q biểu diễn: $75 - 15 = 60$.
}
\end{ex}

\begin{ex}
Số nào có số liền sau là 80?
\shortans{79}
\loigiai{
Số đó là: $80 - 1 = 79$.
}
\end{ex}

\begin{ex}
Số nào có số liền trước là 63?
\shortans{64}
\loigiai{
Số đó là: $63 + 1 = 64$.
}
\end{ex}

\begin{ex}
Có bao nhiêu số tự nhiên lớn hơn 15 và không vượt quá 20?
\shortans{5}
\loigiai{
Các số: 16, 17, 18, 19, 20. Có 5 số.
}
\end{ex}

\begin{ex}
Tổng của số liền trước số 10 và số liền sau số 10 là bao nhiêu?
\shortans{20}
\loigiai{
Tổng: $9 + 11 = 20$.
}
\end{ex}
